\chapter{Related Work}
\label{cha:related-work}

\section{Remote Sensing Image Classification}

Deep learning has become the dominant approach for remote sensing image classification, 
leveraging convolutional neural networks (\gls{cnn}) to extract spatial and spectral features 
from aerial and satellite imagery. While significant progress has been made in classifying 
land-use and scene categories, fire risk classification remains a challenging task due 
to its dependence on vegetation conditions, environmental variations, and fine-grained 
differences between risk levels.


\section{CNN-Based Methods for Remote Sensing Classification}

The article \textit{Remote Sensing Image Scene Classification Meets Deep Learning} 
\cite{cheng2020remote} provides a systematic review of deep learning methods for remote 
sensing image classification, covering more than 160 papers. It compares the accuracies 
of several models using benchmark datasets, including UC-MERCED \cite{yang2010bag}, AID 
\cite{xia2017aid}, and NWPU-RESISC45 \cite{cheng2017remote}.

The UC-Merced dataset contains 2100 images across 21 land-use classes, with a resolution of 
256x256 pixels at 0.3 m spatial resolution. The AID dataset expands on this, featuring 10,000 
images from 30 scene classes, varying in resolution from 8 m to 0.5 m per pixel. NWPU-RESISC45 
is the largest, with 31,500 images spanning 45 categories, covering over 100 countries, with
pixel resolutions ranging from 30 m to 0.2 m.

Although these datasets focus on general land-use and scene classification, they share similarities 
with fire risk classification, particularly in dealing with heterogeneous environments and high 
interclass similarity. In both cases, distinguishing between classes can be challenging due to 
factors such as subtle differences in vegetation density, variations in terrain, and seasonal 
changes affecting the imagery.

\subsection{Deep Learning-Based Scene Classification Algorithms}

Scene classification using deep learning has gained significant attention in remote sensing research. 
The paper \cite{cheng2020remote} categorizes deep learning approaches into three main types: autoencoder-based methods, 
CNN-based methods, and GAN-based methods. While all three have been explored for remote sensing applications, 
CNN-based approaches have demonstrated the highest performance and practical applicability.

\subsubsection{Autoencoder-Based Methods}
Autoencoders have been widely used in earlier years for remote sensing classification due to their ability 
to automatically learn mid-level feature representations from unlabeled data. However, their performance 
is generally inferior to CNN-based approaches, as they do not fully leverage scene class information 
for optimal feature discrimination. Given these limitations, autoencoder-based methods 
are not considered in this study.

\subsubsection{CNN-Based Methods}
CNN-based approaches have become the dominant technique for remote sensing image classification, offering 
state-of-the-art accuracy across various datasets. These methods can be divided into three categories:

\begin{itemize}
    \item \textbf{Pretrained \gls{cnn} as Feature Extractors:} \gls{cnn} pretrained on large datasets such as ImageNet are used to extract 
    deep features from remote sensing images, which are then fed into a separate classifier (e.g., SVM).
    \item \textbf{Fine-Tuning Pretrained \gls{cnn}:} A CNN pretrained on ImageNet is fine-tuned on the target dataset, adjusting its 
    learned features for improved performance.
    \item \textbf{Training \gls{cnn} from Scratch:} A CNN is trained from the ground up using a custom dataset, requiring 
    significantly more labeled samples and computational resources.
\end{itemize}

Among these, fine-tuning pretrained \gls{cnn} has shown the best trade-off between accuracy and computational efficiency, 
particularly for remote sensing applications with limited labeled data. Fine-tuning approaches such as SF-CNN 
(Scale-Free CNN) \cite{SF-CNN} and GBN (Gated Bidirectional Network) \cite{GBN} have achieved accuracies above 98\% 
on benchmark datasets.

For larger-scale and more complex classification tasks, ensemble models like Hydra \cite{hydra} have been
shown to outperform single \gls{cnn}. Hydra employs an ensemble of ResNet and DenseNet architectures,
fine-tuning multiple \gls{cnn} separately and combining their outputs through a majority voting strategy.
An image is classified into the majority class if at least half of the models agree; otherwise, it is labeled as unknown.
This approach allows for improved 
generalization and robustness, particularly in datasets with high interclass similarity.

\paragraph{Hydra Framework}
This framework follows an ensemble-based approach designed to improve classification accuracy while reducing training time. 
The process consists of two main stages:
\begin{enumerate}
    \item An initial CNN is trained and coarsely optimized, forming the "Hydra's body."
    \item The obtained weights are then fine-tuned multiple times using different augmentation techniques, region crop styles, 
    and class weighting strategies to create the ensemble’s individual models, referred to as the "Hydra's heads."
\end{enumerate}

By leveraging this strategy, the ensemble encourages diversity among the individual models, which improves the robustness of predictions. 
A visualization of the Hydra framework is presented in Figure \ref{fig:hydra-testing}.

\begin{figure}[ht]
    \centering
    \includegraphics[width=0.8\textwidth]{../Figures/hydra-testing.png}
    \caption{Hydra framework testing flowchart from \cite{hydra}.}
    \label{fig:hydra-testing}
\end{figure}

The training procedure ensures that each head of the Hydra converges to a different local optimum by employing various 
augmentation techniques such as online data augmentation, different crop strategies, and metadata inclusion. This ensemble 
approach mitigates the limitations of using a single CNN and has demonstrated superior classification performance in remote sensing applications.

At inference time, an input image is passed through all heads of the Hydra, and each head provides a classification score. 
The final prediction is determined through majority voting, which enhances the ensemble’s robustness against individual model biases.

\paragraph{Limitations of Fine-Tuning Pretrained \gls{cnn}}
Despite their effectiveness, fine-tuning pretrained \gls{cnn} introduces certain challenges. Pretrained models are typically optimized for generic 
image classification tasks (e.g., natural images from ImageNet), which may not fully align with the spectral and spatial characteristics of 
remote sensing images. Furthermore, pretrained models may introduce biases based on the original training data, necessitating careful 
selection of hyperparameters and augmentation techniques.

\subsubsection{GAN-Based Methods}
Generative Adversarial Networks (GANs) have also been explored as an alternative approach to remote sensing classification, 
offering the ability to generate synthetic samples for training deep learning models \cite{attgan, martagan}.  
While GAN-based techniques can help mitigate data scarcity, their application in scene classification remains limited.

Compared to CNN-based methods, GANs face two primary drawbacks: (1) they typically require additional classifiers 
for training, making them less efficient as stand-alone classification models, and (2) their overall classification 
accuracy tends to be lower than that of CNN-based approaches. Given these constraints, GAN-based methods are not 
considered in this study.

\section{Conclusion}

Deep learning has significantly advanced remote sensing image classification, with CNN-based 
approaches achieving state-of-the-art performance. Among the different strategies explored 
in the literature, fine-tuning pretrained CNNs has proven to be the most effective for tasks 
with limited labeled data, offering a strong balance between computational efficiency and accuracy.

Ensemble learning methods, such as the Hydra framework, have further demonstrated their 
effectiveness in handling complex classification problems. By leveraging multiple CNN 
architectures and combining their outputs through majority voting, ensemble models enhance 
generalization and reduce the impact of individual model biases. This approach is 
particularly relevant to fire risk classification, where subtle differences between 
risk categories and interclass similarity can pose significant challenges.

Given these findings, this study will explore the application of an ensemble-based 
method for fire risk classification. By integrating models such as ResNet-50, DenseNet-161, 
and EfficientNet-B5, we aim to leverage the strengths of each architecture to improve 
prediction robustness. A majority voting strategy will be employed to ensure stable 
and reliable classifications, mitigating potential biases from individual models. 
